\subsection{Experiment~17: One Threshold for Every Protein}
\label{sec:exp17-calibration-wall}

Every arm of this system emits the same way. Scores are produced per candidate,
a single threshold is chosen to maximise the measure over the whole panel, and
every protein is cut at that threshold. The threshold is fitted; the fact that
one threshold serves all of them is not fitted at all. It is a structural
choice that no experiment in this thesis has examined.

This experiment asks what that choice costs. The answer is larger than any
other margin measured here, and most of it is still unclaimed.

\subsubsection{The Ceiling}

Hold the ordering inside each protein completely fixed and let each protein
choose its own threshold. Nothing about the scores changes and no candidate is
added: the only thing that varies is where each protein's list is cut. The
oracle reads the answer to make that choice, so what follows is a ceiling and
not an opportunity, and it is reported as one throughout.

\begin{table}[htbp]
\centering
\caption{What a per-protein threshold would gain, no-knowledge cells, additions
only, scored over the whole term field.}
\label{tab:exp17-ceiling}
\begin{tabular}{lrrrrrrr}
\toprule
 & & \multicolumn{3}{c}{pooled, IA-weighted} & \multicolumn{3}{c}{per-protein mean} \\
\cmidrule(lr){3-5}\cmidrule(lr){6-8}
Panel & $n$ & base & oracle & gain & base & oracle & gain \\
\midrule
BPO & 1509 & 0.1489 & 0.3631 & +0.2142 & 0.2682 & 0.4867 & +0.2185 \\
MFO & 1129 & 0.3004 & 0.5653 & +0.2649 & 0.4902 & 0.7317 & +0.2415 \\
CCO & 1116 & 0.2855 & 0.5581 & +0.2726 & 0.5228 & 0.7246 & +0.2018 \\
\bottomrule
\end{tabular}
\end{table}

Two properties of this table matter more than its size.

It is present in both statistics and at the same magnitude. The two disagree by
at most 0.07 on any panel, and neither is a rounding of the other: one pools a
confusion matrix weighted by information accretion, the other averages an
$F$-measure over proteins. Almost every promising signal examined in this work
has separated well and then failed to convert between those two. This one does
not.

And the base reproduces the deployed system. Against the deployed arm measured
independently, the per-protein-mean bases here differ by $+0.017$, $+0.004$ and
$+0.003$. Two of the three panels agree to within four thousandths, which is
what establishes that this is the same object and not a favourable
reconstruction of it.

\subsubsection{Where the Margin Lives}

The margin is not spread across the population. At the optimal global
threshold, the tenth of proteins that receive the most predictions receive an
order of magnitude more than the median:

\begin{table}[htbp]
\centering
\caption{The median protein and the top decile by emission, at one global
threshold.}
\label{tab:exp17-tail}
\begin{tabular}{lrrrrr}
\toprule
Panel and group & $n$ & emitted & precision & recall & true terms \\
\midrule
BPO, below p90 & 1357 & 10 & 0.1643 & 0.2859 & 14 \\
BPO, top decile & 152 & 182 & 0.0648 & 0.7464 & 13 \\
CCO, below p90 & 993 & 8 & 0.4012 & 0.5153 & 8 \\
CCO, top decile & 123 & 41 & 0.1713 & 0.8486 & 7 \\
MFO, below p90 & 1015 & 4 & 0.3808 & 0.4740 & 5 \\
MFO, top decile & 114 & 31 & 0.2520 & 0.8570 & 7 \\
\bottomrule
\end{tabular}
\end{table}

The last column is the one to read first. The tail proteins do not have more
true annotations than the median ones: thirteen against fourteen, seven against
eight, seven against five. They are not proteins that do more. They are
proteins to which the system emits eighteen times more.

And their recall is two to three times higher, at a third of the precision. So
the tail is not noise to be swept away. It is where the method finds most of
the answer and buries it. Cutting it back does not remove rubbish; it trades
recall the protein has in abundance for precision it lacks.

That reframes what a per-protein threshold does. It does not calibrate volume
to what each protein deserves. It rescues the proteins that one global
threshold serves worst, and those proteins are identifiable by how much they
emit rather than by anything about their truth.

\subsubsection{A Realisable Rule, and What It Is Not}

The rule tested is the crudest one that can express a per-protein cut: keep the
terms scoring within a fixed fraction $q$ of the protein's own best score. One
scalar, fitted on half the proteins and scored on the other half, in both
directions.

\paragraph{It works, and not by its own name.} Against a global threshold it
gains $+0.022$ to $+0.069$ on the per-protein mean, positive in nine of nine
panel-by-artefact combinations with intervals excluding zero. Giving a protein
some \emph{other} protein's peak and refitting $q$ kills it: negative in
eighteen of eighteen.

But the peak does not predict how many terms are true. The rank correlation
between the two runs from $+0.03$ to $+0.11$ and changes sign between
artefacts, so it is zero and not even stably signed. The rule is therefore not
calibrating volume to evidence, whatever its description suggests. What it
does is keep the terms within a relative band of the top, which is an operation
on the shape of a protein's score distribution and not on its level. A protein
whose scores fall away sharply keeps few; one with a flat summit keeps many.

\paragraph{And it is not abstention.} By construction the rule cannot silence
anyone: the cut always lands where the protein is still emitting, at every $q$
including one. It is a trimming rule. The oracle's silence, which reaches a
third of the panel, is a separate lever and remains without a key.

\paragraph{The gain is where the mechanism says it should be.} Fitting on the
whole population and evaluating the difference separately, the tail gains four
to six times what the median does, $+0.105$ to $+0.235$ against $+0.022$ to
$+0.078$, with eight of nine tail cells excluding zero. So $+0.03$ over the
whole panel is a large improvement to a tenth of it, not a small improvement to
everyone.

\subsubsection{The Two Measures Disagree About What to Deploy}

Against a global threshold the rule wins on both statistics. Against a fixed
top-$k$, a baseline that emits the same number of terms for every protein and
costs nothing to implement, they part company.

\begin{table}[htbp]
\centering
\caption{Per-protein rescaling against a fixed top-$k$, cross-fitted, over
three artefacts.}
\label{tab:exp17-metrics}
\begin{tabular}{lll}
\toprule
Statistic & Rule applied to all & Rule applied to the top quintile \\
\midrule
Per-protein mean & wins, median $+0.0136$ & falls to a third \\
Pooled IA-weighted & \textbf{loses} to fixed top-$k$ & wins, $+0.015$ to $+0.047$ \\
\bottomrule
\end{tabular}
\end{table}

Under the pooled weighted measure the unrestricted rule actively harms the
median protein, which is what makes it unstable there, and restricting it to
the highest-emitting fifth removes the harm and turns eight of nine panels
positive. Under the per-protein mean the median also gains, so restricting the
rule throws that gain away.

The targeting is not merely doing less. Applying the rule to a randomly drawn
decile of the same size, over twenty draws per cell, is beaten by volume-based
targeting in nine of nine panels under the pooled measure, by a factor of three
to four. Under the per-protein mean the random decile matches it in six of
nine, so there the restriction really is buying by doing less, and the criterion
is decoration.

\textbf{The finding is that pair of rows, not either rule.} A per-protein mean
rewards repairing the few proteins served worst, because each protein counts
once. A pooled sum weighted by information accretion penalises disturbing the
many that were already fine, because what is trimmed carries mass. The two
answers are not reconcilable by fitting, and the system reports to the
community benchmark under one measure while its own surfaces read the other.
Which rule to deploy therefore has no metric-free answer, and choosing is a
decision to be taken openly rather than settled inside a hyperparameter.

\subsubsection{What Bounds It}

The gains are small in absolute terms and not small against the bases they
improve: $+0.032$ on $0.1452$ is twenty-two per cent, $+0.032$ on $0.1625$ is
twenty, $+0.026$ on $0.1438$ is eighteen. Both readings are true and they
persuade different readers, so both are given.

The oracle reads the answer. The ceiling of $+0.20$ to $+0.27$ sizes the
question and not the opportunity, and the realisable rule recovers between a
ninth and a third of it. What claims the rest is unknown.

Two panels cannot be measured. The per-protein artefact refuses to write the
prior-knowledge cells for two of three aspects, because the evaluation kernel
pools its confusion matrix over rows it excludes from the population it
normalises by. That refusal is correct and the defect it names is real, but it
leaves the largest cells of the frame outside this experiment.

The directed rule carries two fitted parameters, not one, and the threshold it
uses is discretised: a protein's peak takes sixty distinct values on a grid of
ninety-nine, so a continuous $q$ is in practice sixty steps. And every number
here is one temporal window and one reference bank, so the stationarity of the
annotation process, which is not established anywhere in this work, bounds how
far any of it carries forward.
